\section{Fractional cardinality and spatial certificates}
\label{sec:cardinality-spatial}
The following optimization problems have explicit solutions. Their role is
to distinguish the size of certificates produced by specified spatial
oracles from the difficulty of finding an optimizer. The obstruction uses a
concave quadratic penalty and an affine balance; it is a separate
construction from the positive multilinear gap examples above.

Fix integers $n\ge2$ and $0\le k\le n-1$, and write
\begin{equation}\label{eq:cardinality-problem}
 K=k+\tfrac12,\qquad
 \mathcal F_{n,k}=\{x\in[0,1]^n:\textstyle\sum_i x_i=K\},\qquad
 F(x)=\sum_{i=1}^n x_i(1-x_i).
\end{equation}
The dimension parameter $n$ counts original coordinates. A vertex of
$\mathcal F_{n,k}$ has exactly $k$ coordinates equal to one, one equal to
$1/2$, and the others zero: two strictly fractional coordinates would admit
a small opposite perturbation, whereas a point with at most one fractional
coordinate is determined by its active bounds and the balance. Concavity
and compactness therefore give
\begin{equation}\label{eq:cardinality-optimum}
 \min_{\mathcal F_{n,k}}F=\tfrac14.
\end{equation}
For each oracle below, a certified region is either infeasible or has
oracle bound at least $T_*$. A certificate is a finite cover of
$\mathcal F_{n,k}$ by such regions, as in
Section~\ref{subsec:certificate-model}. For absolute tolerance
$\epsilon\ge0$ we grant the best incumbent and set
$T_*=1/4-\epsilon$. Pruning at $U-\epsilon$, for any feasible incumbent
$U\ge1/4$, implies this condition. Infeasible regions contain no witness
and do not help cover the feasible set. All lower bounds also hold for a
weaker valid oracle: certifying a region with that oracle certifies it
with the stronger oracle used in the proof.

There are close precedents in other certificate models. Jarre's binary
knapsack construction retains an exponential branch-and-bound obstruction
with a max-cut semidefinite bound \citep[Sections 2--3]{Jarre2018}.
Coniglio's hyperplane-clustering analysis gives a related spatial
obstruction under midpoint branching on symmetric coordinate domains
\citep[Assumption 1, Propositions 2--4, Appendix C in the review version]{Coniglio2026}.
Our arguments concern arbitrary coordinate restrictions and a fixed
objective gap under the explicitly defined node oracles. We make no
priority claim for exponential spatial lower bounds in general. The
fractional-cardinality moments used later are classical; the additional
argument places them in spatial regions and counts how many witnesses a
certified region can contain.

\subsection{Chords and endpoint avoidance}
For a box $B=\prod_i[a_i,b_i]\subseteq[0,1]^n$, including degenerate
intervals, define
\begin{equation}\label{eq:cardinality-chord}
 \ell_i(x_i)=a_ib_i+(1-a_i-b_i)x_i,\qquad
 x_i(1-x_i)-\ell_i(x_i)=(x_i-a_i)(b_i-x_i).
\end{equation}
Its chord bound is
\[
 \operatorname{LB}_{\rm ch}(B)=
 \min_{x\in B\cap\mathcal F_{n,k}}\sum_i\ell_i(x_i),
\]
with $+\infty$ for an empty feasible intersection. If $g_i$ is any convex
underestimator of $x_i(1-x_i)$ on $[a_i,b_i]$, convexity between the
endpoints gives $g_i\le\ell_i$. This also holds on a singleton interval.
Thus the infimum for any separable convex underestimator $\sum_i g_i$
is at most $\operatorname{LB}_{\rm ch}(B)$; an infimum is used because
arbitrary convex underestimators need not attain their infima on a closed
interval without lower semicontinuity at its endpoints.

For clarity, two common formulations yield exactly these chords. For
$y_i=1-x_i\in[1-b_i,1-a_i]$, both McCormick lower planes for
$w_i=x_iy_i$ become $w_i\ge\ell_i(x_i)$. The upper planes remain
compatible: the exact product lies above the chord and below each upper
plane. The separable $\alpha$BB expression
$x_i(1-x_i)-\alpha(x_i-a_i)(b_i-x_i)$ has second derivative
$2(\alpha-1)$, so the smallest convexifying value $\alpha=1$ gives the
chord. Larger valid values give weaker underestimators. These statements
concern the specified separable formulations, without additional coupled
cuts or nonlinear subproblem solution rules.

We isolate the counting argument so that later oracle changes do not
obscure it. Let $n=k+m+z$, where $k,z\ge0$ and $m\ge1$ are integers.
For each ordered partition $(H,M,Z)$ of $[n]$ with these sizes, put
\begin{equation}\label{eq:cardinality-witness}
 p=\frac1{2m},\qquad w=\ind_H+p\ind_M.
\end{equation}
These $n!/(k!m!z!)$ distinct witnesses belong to $\mathcal F_{n,k}$.
For an arbitrary product $S=\prod_i S_i$, with nonempty
$S_i\subseteq[0,1]$, define
\begin{equation}\label{eq:endpoint-exclusions}
 A(S)=\{i:0\notin S_i\},\quad D(S)=\{i:1\notin S_i\},\quad
 R(S)=A(S)\cup D(S).
\end{equation}
For boxes these are $A=\{i:a_i>0\}$ and $D=\{i:b_i<1\}$.

\begin{lemma}[Endpoint-avoidance count]\label{lem:endpoint-count}
Suppose each certified product domain that contains a witness has
$|R(S)|\ge q$, where $q\ge0$ is real. Any certified cover of
$\mathcal F_{n,k}$ contains at least
\begin{equation}\label{eq:cardinality-cover}
 \mathcal N(q):=
 \min\left\{\left(\frac n{n-k}\right)^{q/2},
                 \left(\frac n{n-z}\right)^{q/2}\right\}
\end{equation}
regions. The same conclusion follows from the stronger hypothesis
$|R(S)|>q$. Both denominators are positive, and $\mathcal N(0)=1$.
\end{lemma}
\begin{proof}
Choose the ordered partition uniformly. A contained witness must have
$Z\cap A=\varnothing$ and $H\cap D=\varnothing$. For a fixed set of
$s$ excluded positions, a uniformly chosen $z$-subset avoids it with
probability
\[
 \frac{\binom{n-s}{z}}{\binom nz}
 =\prod_{j=0}^{s-1}\frac{n-z-j}{n-j}
 \le\left(\frac{n-z}{n}\right)^s
 \quad(s\le n-z).
\]
For $s>n-z$ the probability is zero, so the upper bound still holds;
for $s=0$ the product is one. The analogous inequality holds for $H$.
No independence between the two avoidance events is required. A fixed
region's witness fraction is at most
\[
 \min\left\{\left(\frac{n-z}{n}\right)^{|A|},
               \left(\frac{n-k}{n}\right)^{|D|}\right\}.
\]
Since $|A|+|D|\ge |R|\ge q$, at least one exponent is at least $q/2$.
The displayed fraction is consequently at most $1/\mathcal N(q)$.
Regions containing no witness have fraction zero. The union bound over a
cover proves the assertion, whether or not its regions overlap. Bases
equal to one when $k=0$ or $z=0$ cause no exception.
\end{proof}

\begin{theorem}[Separable spatial certificates]\label{thm:chord-cover}
Let $n=k+m+z$, with integers $k,z\ge0$, $m\ge1$, $n\ge2$, and
$0\le\epsilon<1/4$. For problem~\eqref{eq:cardinality-problem}, every
cover by boxes certified by the chord oracle at
$T_*=1/4-\epsilon$ has at least $\mathcal N(h)$ members, where
\begin{equation}\label{eq:chord-h}
 h=m(1/2-2\epsilon).
\end{equation}
This includes arbitrary real coordinate split points and all separable
convex underestimators dominated by the chord oracle. For
$k=m=z=t$, $n=3t$, the bound is
$(3/2)^{t(1/4-\epsilon)}$, hence $(3/2)^{n/24}$ at $\epsilon=1/8$.
\end{theorem}
\begin{proof}
In a box containing $w$, the chord values at its zero and unit witness
coordinates vanish. On $U=[n]\setminus R$ the interval is $[0,1]$,
whose chord vanishes identically. Consequently
\[
 T_*\le\operatorname{LB}_{\rm ch}(B)
 \le\sum_{i\in M\cap R}\ell_i(p)
 \le |M\cap R|p(1-p)<\frac{|M\cap R|}{2m}.
\]
The last strict inequality is applicable because $T_*>0$ forces
$M\cap R\ne\varnothing$. Thus $|R|>2mT_*=h$, and
Lemma~\ref{lem:endpoint-count} applies. A complete coordinate split tree
provides such a cover, regardless of split locations or branching order.
\end{proof}

\subsection{Upper certificates, tightening, and tolerances}
\begin{proposition}[Explicit upper certificates]\label{prop:chord-upper}
For $0<\epsilon<1/4$, a chord certificate for
\eqref{eq:cardinality-problem} has at most $2^{n+2}-3$ nodes and
$2^{n+1}-1$ leaves. For $\epsilon=0$, a midpoint certificate has exactly
$\binom{n+1}{k+1}$ leaves if only the stopping rules in the proof are used.
In particular its node count is
$O(n^{\min\{k+1,n-k\}})$ when that exponent is fixed.
\end{proposition}
\begin{proof}
Choose $0<\alpha<1/2$ with
$\alpha(1-\alpha)=1/4-\epsilon$. Process coordinates in order. At each
surviving node split the next coordinate at $\alpha$, then split its
upper child at $1-\alpha$. The middle interval has constant chord
$\alpha(1-\alpha)$, and all other chords are nonnegative, so the middle
child is certified. Continue on the low and high intervals. At a final
corner box the chord sum is $(1-\alpha)\|x-v\|_1$, where
$v\in\{0,1\}^n$ is its corner. On the demand slice,
$\|x-v\|_1\ge|K-\sum_i v_i|\ge1/2$. This gives a bound
$(1-\alpha)/2>1/4>T_*$. Processing coordinate $i$ adds four nodes per
surviving prefix, of which there are $2^{i-1}$. The stated counts follow.
For a strict pruning rule choose a slightly larger middle-chord target
still below $1/4$.

For exact certification split at $1/2$ instead. If a prefix has $c$ high
intervals and $d$ low intervals, its chord bound is at least each of
\[
 \frac{c-K}{2},\qquad \frac{K-(n-d)}2.
\]
Indeed, the sum on high coordinates is at most $K$, whereas the sum on
low coordinates is at least $K-(n-d)$; all other chord contributions
are nonnegative. Stop at $c=k+1$ or $d=n-k$. No path survives depth $n$.
With remaining high and low stopping counts $a,b$, the number of leaves
satisfies $L(a,b)=L(a-1,b)+L(a,b-1)$ and
$L(0,b)=L(a,0)=1$, hence $L(a,b)=\binom{a+b}{a}$. Start at
$(a,b)=(k+1,n-k)$. A full binary tree has twice as many leaves minus one
nodes. This proves the exact count and polynomial assertion.
\end{proof}

The balanced fixed-tolerance family therefore has certificate size
$2^{\Theta(n)}$ under the chord model. If both $k$ and $n-k$ are
proportional to $n$, choose $m,z$ proportional to $n-k$ to obtain the
same exponential order. Fixed demand level, or fixed complementary demand
level, admits the polynomial certificate just proved. At
$\epsilon\ge1/4$, the root chord bound zero already meets the target;
no positive exponent is asserted there.

The chord error is genuinely quadratic in interval width:
\begin{equation}\label{eq:second-order-chord}
 0\le F(x)-\sum_i\ell_i(x_i)
 \le\frac14\sum_i(b_i-a_i)^2
 \le\frac n4\max_i(b_i-a_i)^2.
\end{equation}
In particular, the difference between the exact node optimum and the
chord bound is at most this quantity, by evaluating $F$ at a chord
minimizer. Thus the fixed-tolerance lower bound persists despite this
second-order error estimate; it does not require a first-order local
approximation error.

\begin{proposition}[Charged reductions and equality tolerance]
\label{prop:charged-tolerance}
Suppose a run keeps feasible points covered by its final boxes together
with discarded boxes certified by the same oracle at $T_*$. If it processes
$N$ nodes and discards at most $a$ such boxes per node, a cover lower bound
$L$ implies $N\ge L/(a+1)$. Pure feasibility reductions need not be
charged when they discard no feasible point.

For a demand feasibility tolerance $0\le\delta<1/2$, replace
$\mathcal F_{n,k}$ by
$\mathcal F^\delta_{n,k}=\{x\in[0,1]^n:|\sum_i x_i-K|\le\delta\}$.
Its exact optimum is $1/4-\delta^2$. If
$\epsilon+\delta^2<1/4$, a chord certificate at absolute tolerance
$\epsilon$ for this enlarged set satisfies Theorem~\ref{thm:chord-cover}
with $\epsilon$ replaced by $\epsilon+\delta^2$.
\end{proposition}
\begin{proof}
The augmented cover has at most $(a+1)N$ members. A discarded open
coordinate slab may be charged by its closed box only when that box is
certified; this condition is part of the assertion, not an unproved
closure assumption. Subsequent product-domain results can charge the
removed coordinate set itself.

For a fixed sum $s\in[k+1/2-\delta,k+1/2+\delta]\subset(k,k+1)$,
the same vertex argument as above gives minimum penalty
$(s-k)(1-s+k)$. Its minimum on this interval is
$1/4-\delta^2$, attained by $k$ ones and a remaining coordinate
$1/2-\delta$ or $1/2+\delta$. A tolerance-feasible incumbent is therefore
at least this value. The exact-balance witnesses remain feasible for the
enlarged problem, and the chord bound over the enlarged feasible set is
at most its value at each witness. The proof of
Theorem~\ref{thm:chord-cover} applies at the positive target
$1/4-(\epsilon+\delta^2)$.
\end{proof}
The restriction $\delta<1/2$ is material: once the tolerance interval
contains an integer sum, the optimum is zero. An unrestricted formula
$1/4-\delta^2$ would then be false. Relative pruning in the convention
$\operatorname{LB}\ge(1-\theta)U$, $0\le\theta<1$, on the original
problem implies the absolute target with $\epsilon=\theta/4$.

\subsection{Full second moments and box RLT}
For $B=\prod_i[a_i,b_i]\subseteq[0,1]^n$, let
$\operatorname{LB}_{\rm SDP}(B)$ be the infimum of
$\sum_i(\mu_i-X_{ii})$ over symmetric $X$ and $\mu\in\R^n$ satisfying
\begin{align}
 a\le\mu\le b,\quad &\sum_i\mu_i=K,\quad
 \begin{pmatrix}1&\mu^T\\\mu&X\end{pmatrix}\succeq0,
 &\sum_jX_{ij}=K\mu_i\quad(i\in[n]),\label{eq:sdp-model}
\end{align}
and all pairwise box inequalities
\begin{align}
 X_{ij}-a_i\mu_j-a_j\mu_i+a_ia_j&\ge0,\notag\\
 b_j\mu_i-X_{ij}-a_ib_j+a_i\mu_j&\ge0,\notag\\
 b_i\mu_j-X_{ij}-b_ia_j+a_j\mu_i&\ge0,\notag\\
 b_ib_j-b_i\mu_j-b_j\mu_i+X_{ij}&\ge0
 \qquad(i,j\in[n]).\label{eq:box-rlt}
\end{align}
Indices may coincide. These are the linearizations of all products of
$x_i-a_i$ and $b_i-x_i$. Exact evaluation $(\mu,X)=(x,xx^T)$ shows
validity. As before the infimum of an empty relaxation is $+\infty$,
and a region is certified when its bound is at least
$T_*=1/4-\epsilon$ (or it is infeasible).

\begin{lemma}[Restricted fractional second moments]\label{lem:sdp-witness}
Let $k,m,z\ge1$, and let $B$ contain a witness~\eqref{eq:cardinality-witness}.
If $|R(B)|<\min\{k,z\}$, the model
\eqref{eq:sdp-model}--\eqref{eq:box-rlt} has a feasible point of value
$|M\cap R|p(1-p)$.
\end{lemma}
\begin{proof}
Fix $x_i=w_i$ deterministically on $R$ and put
\[
 U=[n]\setminus R,\quad s=|U|,\quad
 t=K-\sum_{i\in R}w_i=|H\cap U|+p|M\cap U|.
\]
At least one unit and one zero witness coordinate remain, so $s\ge2$
and $1\le t\le s-1$. On $U$ set
\[
 c=t/s,\quad d=t(t-1)/(s(s-1)),\quad
 \mu_i=c,\quad X_{ii}=c,\quad X_{ij}=d\ (i\ne j).
\]
On $R$ set $\mu_i=w_i$, and set $X_{ij}=\mu_i\mu_j$ whenever at least
one index lies in $R$. First moments meet the box bounds because the
intervals on $U$ are $[0,1]$ and $w\in B$. Their sum is $K$.
The covariance is zero on rows and columns in $R$, and on $U$ it is
\[
 X_{UU}-\mu_U\mu_U^T
 =(c-d)(I-J/s),\qquad c-d=\frac{t(s-t)}{s(s-1)}\ge0.
\]
It is positive semidefinite, proving augmented PSD by the Schur complement
of the entry one. For $i\in U$,
$c+(s-1)d=tc$ gives $\sum_jX_{ij}=Kc$; for $i\in R$, the equality
is $\sum_jX_{ij}=w_i\sum_j\mu_j=Kw_i$.

For distinct $i,j\in U$, the four slack-product expectations are
\[
 d,\quad c-d,\quad c-d,\quad
 1-2c+d=\frac{(s-t)(s-t-1)}{s(s-1)},
\]
all nonnegative. For $i=j\in U$ they are $c,0,0,1-c$.
If at least one index is in $R$, the expectation factors into a
nonnegative deterministic slack and a nonnegative expected slack; this
also covers repeated restricted indices and degenerate intervals.
Finally $\mu_i-X_{ii}=0$ on $U$, and on $R$ it is $w_i(1-w_i)$.
\end{proof}

\begin{theorem}[SDP--RLT cover bound]\label{thm:sdp-cover}
Let $n=k+m+z$, where $k,m,z\ge1$ are integers, and let
$0\le\epsilon<1/4$. Every cover for
\eqref{eq:cardinality-problem} certified by
\eqref{eq:sdp-model}--\eqref{eq:box-rlt} at $T_*=1/4-\epsilon$ has
at least $\mathcal N(q)$ members, where
$q=\min\{k,z,m(1/2-2\epsilon)\}$.
In the balanced case $k=m=z=t$, this is the chord lower bound in
Theorem~\ref{thm:chord-cover}.
\end{theorem}
\begin{proof}
If a containing box has $|R|<q$, Lemma~\ref{lem:sdp-witness} gives a
feasible objective at most $|R|p(1-p)<q/(2m)\le T_*$, where the
strict comparison holds also when $R=\varnothing$ because $q>0$.
The box cannot be certified. Therefore $|R|\ge q$ on every certified
box containing a witness, and Lemma~\ref{lem:endpoint-count} applies.
The diagonal mixed inequality gives
$X_{ii}\le(a_i+b_i)\mu_i-a_ib_i$, so the SDP bound dominates the
chord bound. The upper certificates of Proposition~\ref{prop:chord-upper}
remain valid.
\end{proof}
This resolves the behavior under spatial branching for this SDP--RLT
model; it is not merely a calculation of its root value.

One can also grant all products of two affine inequalities valid on
$B\cap\mathcal F_{n,k}$ without strengthening this oracle. Indeed,
linear-programming duality represents any such nonnegative affine
function, on a nonempty polytope, as a nonnegative constant plus a
nonnegative combination of box slacks plus a multiple of
$e(x)=\sum_i x_i-K$. The representation is an identity of affine
polynomials, including for degenerate boxes. Expanding the product of
two representations leaves only existing nonnegative slack moments and
terms $L[e v]=0$ with $v$ affine. This assertion concerns affine
inequalities in the original $x$ coordinates. It does not grant all
valid affine cuts in the lifted $(\mu,X)$ coordinates.
